\title{Byte Latent Transformer: Patches Scale Better Than Tokens}

\begin{abstract}
We propose the Byte Latent Transformer ({\textsc BLT}), a novel LLM architecture that operates directly at the byte level and, for the first time, achieves parity with tokenization-dependent LLMs on large-scale tasks, while delivering notable gains in inference speed and resilience. In {\textsc BLT}, raw bytes are encoded into adaptively sized patches, which become the main computational units for the model. These patches are formed by analyzing the entropy of the upcoming byte, allowing the allocation of greater computational resources and representational power to more complex data segments. We conduct the first flop-budgeted scaling analysis for byte-level models, extending up to 8B parameters and 4T bytes processed during training. Our findings validate that it is possible to effectively scale byte-oriented models without relying on static vocabularies. Efficiency during both training and inference benefits from the dynamic creation of longer patches in low-complexity regions, along with observed qualitative advancements in tasks involving reasoning and coverage of rare cases. Ultimately, given equal inference resources, {\textsc BLT} achieves superior scaling compared to approaches based on tokenization, by jointly increasing patch length and model capacity.
\end{abstract}

\section{Introduction}
\label{section:intro}

We present the Byte Latent Transformer~(\textbf{BLT}), an architecture that eliminates the need for traditional tokenization by directly operating on raw byte sequences. Notably, this approach attains parity with tokenization-based models in large-scale settings while achieving notable gains in computational efficiency and resilience (\S\ref{sec:robustness}). In contrast, contemporary large language models (LLMs) employ a nearly end-to-end training strategy with the exception of tokenization—a heuristic preprocessing stage that segments byte streams into a fixed vocabulary of tokens. This step introduces biases into the compression of input strings and gives rise to limitations, including sensitivity to specific domains or modalities~\citep{dagangetting}, vulnerability to noise in input data~(\S\ref{sec:robustness}), deficits in orthographic representation~\citep{edman2024cute}, and disparities in multilingual capabilities~\citep{liang2023xlm,petrov2024language,limisiewicz2024myte}.

Tokenization has historically been crucial, as training large language models on raw bytes without it proves extremely expensive at scale because of the extended sequence lengths involved~\citep{xue2022byt5}. Earlier studies have addressed this challenge by designing more efficient self-attention mechanisms~\citep{el2020characterbert, clark2022canine} or by adopting architectures that avoid attention entirely~\citep{wang2024mambabyte}~(\S\ref{section:related_work}). Nonetheless, these strategies primarily facilitate the training of \textit{small models}. For large-scale models, the bulk of computational overhead stems from the sizeable feed-forward network layers applied to each byte, rather than from the attention component.

To optimize computational resource allocation, we introduce an adaptable, learnable strategy for aggregating bytes into \textit{patches}~(\S\ref{section:patching}), along with a novel model design that integrates both byte-level and patch-level signals. In contrast to tokenization, {\textsc BLT} does not rely on a predetermined vocabulary for patches. Instead, groups of bytes—regardless of their configuration—are projected to latent patch embeddings using efficient, trainable encoder and decoder components. Our experiments demonstrate that this methodology allows for \textit{greater} compute efficiency compared to models dependent on tokenization.

Traditional tokenization-based LLMs dedicate an equal amount of computational resources to each token, resulting in a compromise between computational efficiency and predictive performance. This is due to the use of compression heuristics in token formation, which often do not correspond to the actual difficulty of making accurate predictions. A key principle of our proposed architecture is adaptive computation—assigning greater resources to parts of the input that require it. For instance, deploying a large transformer block for predicting the endings of most words is unnecessary, as these predictions are generally straightforward and entail low uncertainty, in stark contrast to the more complex task of generating a sentence’s initial word. {\textsc BLT} incorporates this philosophy in its design~(\S\ref{section:architecture}), employing a structure of three transformer blocks: two compact byte-level \textit{local models} paired with a more substantial, global \textit{latent transformer}~(\autoref{fig:arch_high_level}). 

In order to effectively group bytes into patches and judiciously manage computation, {\textsc BLT} partitions data according to the entropy associated with predicting the subsequent byte. This strategy yields contextual segments of bytes with roughly consistent information density, ensuring computation is distributed where prediction complexity is highest.

We introduce the first study to systematically scale byte-level models under a fixed flop budget, extending up to 8 billion parameters and 4 trillion bytes of training data. This work demonstrates that it is possible to perform large-scale, end-to-end model training directly from bytes, without relying on preset token vocabularies. In terms of training performance under controlled computational expenditure,\footnote{The computational expenditure of a model is assessed based on the tally of Floating Point OPerations (flops) required.} {\textsc BLT} performs comparably to Llama 3 while reducing inference flops by as much as 50\%~(\S\ref{section:scaling}).

Additionally, we provide evidence that modeling directly from raw bytes yields clear benefits in capturing infrequent patterns present in the data. Models trained using {\textsc BLT} demonstrate greater robustness to noisy inputs when compared to tokenizer-based counterparts, and they also show improved performance in character-level linguistic tasks such as orthographic and phonological processing, as well as low-resource machine translation~(\S\ref{sec:robustness}).

Moreover, employing {\textsc BLT} models allows simultaneous increases in both model size and patch size without exceeding the preset inference flop limit. Utilizing larger patch sizes generally leads to computation savings, which in turn can be used to increase the scale of the global latent transformer, as it is invoked less frequently. Our scaling investigations with a fixed inference-flop framework~(\autoref{fig:fixed_inference_scaling}) reveal that these approaches exhibit far superior scaling behaviors in comparison to tokenization-dependent architectures.

In conclusion, our work offers the following key contributions: 1) We present {\textsc BLT}, an architecture for byte latent llms that employs dynamic compute allocation to enhance flop efficiency, 2) Our results demonstrate flop-controlled training parity with Llama 3 up to the 8B parameter regime, with the added flexibility to exchange slight decreases in evaluation metrics for as much as 50\% improvement in flop efficiency, 3) {\textsc BLT} models open a novel scaling approach for llms, enabling increases in model capacity without exceeding a predetermined inference budget, 4) We provide evidence that {\textsc BLT} models exhibit superior resilience to input noise and possess a heightened sensitivity to sub-word features in input sequences, which are frequently overlooked by token-based llms. The complete training and inference implementation for {\textsc BLT} is publicly available at~\url{https://github.com/facebookresearch/blt}.

\section{Patching: From Individual Bytes to Groups of Bytes}
\label{section:patching}

Dividing a byte sequence into \textit{patches} enables {\textsc BLT} to flexibly assign computational resources according to contextual requirements. Figure~\ref{fig:patching_types} presents a variety of approaches for organizing bytes into patches. More precisely, a patching function $f_p$ partitions a byte sequence $\pmb{x} = \{x_i\,|\,i=1,\ldots,n\}$ of length $n$ into a collection of $m < n$ patches $\pmb{p} = \{p_j\,|\,j=1,\ldots,m\}$ by assigning each $x_i$ a value in \{0,1\}, with 1 denoting the initiation of a new patch. For both token-based architectures and patch-based approaches, the computational load when processing data primarily hinges on the number of forward steps taken by the core Transformer component. In {\textsc BLT}, this number corresponds to the quantity of patches derived from the chosen patching strategy. Thus, the typical patch length, or \textit{patch size}, plays a central role in determining processing costs at both training and inference stages for a specified patching scheme~(\S\ref{section:flops}).

The following sections present three specific patching methods: partitioning by a constant number of bytes per patch~(\S\ref{section:static-patch}), dividing on whitespaces~(\S\ref{section:white-patch}), and employing entropy-based segmentation using a lightweight byte language model~(\S\ref{section:dyn-patch}). Subsequently, we elaborate on the idea of incremental patching and distinguish patching from conventional tokenization procedures~(\S\ref{section:bpe-patch}).

\subsection{Strided Patching Every K Bytes}
\label{section:static-patch}

One of the simplest approaches for grouping bytes involves forming fixed-size patches of length $k$, as introduced in MegaByte~\citep{yu2023megabyte}. Utilizing a consistent stride not only simplifies the implementation of both training and inference, but also offers an intuitive way to adjust average patch size, thus allowing for straightforward control over computational cost. Despite these advantages, fixed-size patching presents notable limitations. Primarily, computational resources are not adaptively assigned based on content complexity: transformer step $j$ might be underutilized when handling less informative tokens like whitespace in code, whereas regions with high informational density, such as mathematical expressions, may receive inadequate compute. Furthermore, such rigid partitioning can produce variable and context-insensitive splits, for instance, segmenting identical words differently across occurrences.

\subsection{Space Patching}
\label{section:white-patch}

\citet{slagle2024spacebyte} introduces a straightforward yet powerful modification to the strided patching approach by generating new patches immediately after encountering space-like bytes\footnote{
    Space-like bytes refer to any byte that does not correspond to a Latin letter, digit, or utf-8 continuation byte. Furthermore, each patch is required to include at least one byte that is not space-like. }, which commonly act as natural boundaries for linguistic segments in many languages. The Space patching method assigns an additional latent transformer step (i.e., increased computational effort) to each word, thereby guaranteeing consistent patching of words across various sequences and allocating additional computation to those challenging predictions that frequently follow spaces. For instance, generating the initial byte of the response to ``Who composed the Magic Flute?\uline{\hspace{.6em}}'' is considerably more challenging compared to subsequent bytes after the ``M,'' because the initial letter heavily limits the set of plausible continuations, thus rendering the rest of the answer, such as ``Mozart,'' much easier to generate. Nevertheless, Space patching is not universally suitable for all linguistic contexts or domains and is notably unable to adjust the size of patches. In the following, we present a novel patching scheme informed by the observation that predicting the leading bytes of words tends to be most challenging; moreover, this approach also provides a direct way to modulate patch sizes.

\subsection{Entropy Patching: Using Next-Byte Entropies from a Small Byte LM}
\label{section:dyn-patch}

Instead of employing a rule-based heuristic like whitespace, our methodology adopts a data-centric strategy to locate points of high uncertainty in next-byte predictions. We propose \textit{entropy patching}, a technique that leverages entropy measures to determine the boundaries of patches.

To implement this, we train a compact byte-level auto-regressive language model on the {\textsc BLT} training dataset. Using this model, we calculate the entropy associated with the prediction of the subsequent byte, given the language model’s probability distribution $p_{e}$ across the byte vocabulary $\mathcal{V}$:

\begin{equation}
H(x_i) = \sum_{v \in \mathcal{V}} p_{e}(x_i=v|\pmb{x}_{<i}) \log p_{e}(x_i=v|\pmb{x}_{<i})
\end{equation}

We explore two approaches for determining patch boundaries using the entropies $H(x_i)$. The initial method selects points whose entropy surpasses a fixed global threshold, as depicted in~\autoref{fig:patching}. The alternative method detects points with entropy values that are notably higher than their immediate predecessors. This latter strategy can also be understood as locating positions where the generally decreasing trend of entropy within a patch is disrupted.

\begin{align*}
    \text{Global Constraint}\hspace{1cm}H(x_t)& > \theta_g\\
    \text{Approx. Monotonic Constraint}\hspace{1cm}H(x_t)& - H(x_{t-1}) > \theta_r
\end{align*}

Patch boundaries are determined through an efficient preprocessing stage carried out as part of the data loading process. This approach contrasts with~\citet{nawrot-etal-2023-efficient}, where a classifier is trained to recognize patch boundaries using entropy-based criteria. In our experiments~(\S\ref{section:experiments}), we evaluate and contrast these two techniques for separating low entropy from high entropy bytes.

\subsection{The Byte-Pair Encoding (BPE) Tokenizer and Incremental Patching}
\label{section:incremental-patching}
\label{section:bpe-patch}

Numerous contemporary llms, such as our reference model Llama 3, implement subword tokenizers like bpe~\citep{gage1994new, sennrich-etal-2016-neural}. Throughout this paper, we use the term ``tokens'' to indicate byte segments that are sampled from a \textit{finite} vocabulary fixed ahead of training, whereas ``patches'' denote adaptively formed byte sequences that lack a predetermined vocabulary. A key distinction between these representations is that in the case of tokens, the model cannot access the raw byte-level features directly.

An important advancement that {\textsc BLT} introduces compared to models relying on tokenization is a reconfiguration of the balance between vocabulary size and computational requirements. In traditional large language models, expanding the vocabulary leads to tokens that encapsulate more information, thus lowering the total number of decoding steps required. However, this also results in a larger output layer in the final projection, increasing the model’s parameter count and memory footprint. Due to this trade off, tokenization-based methods have limited flexibility in significantly altering token sizes or reducing inference costs. As an illustrative case, Llama 3 managed to raise its average token length from 3.7 to 4.4 bytes; nonetheless, this gain necessitated scaling up its embedding matrix by a factor of four relative to Llama 2.

During generation, {\textsc BLT} must determine at each position in the byte sequence whether it coincides with a patch boundary, as this informs whether additional computation from the Latent Transformer should be engaged. Notably, this boundary determination must be performed based solely on the available prefix, independent of future sequence elements yet to be produced. Consequently, any patching strategy must avoid dependence on ungenerated bytes for making segmentation decisions. More formally, we say a patching scheme $f_p$ is incrementally applicable if it meets the following criterion:
$$f_p(\pmb{x}_{<i}) = f_p(\pmb{x})_{<i}$$
The bpe algorithm does not satisfy this incremental requirement: tokenizing an identical prefix may yield different outputs contingent upon the following bytes, thus violating the above property\footnote{Incorporating an explicit delimiter token to mark patch boundaries can make bpe behave as an incremental patching method, but at the cost of inflating the sequence length.}.

\section{{\textsc BLT} Architecture}
\label{section:architecture}
{\textsc BLT} consists of a large-scale, globally conditioned autoregressive language model designed to process patch representations, in addition to two lightweight, specialized models responsible for mapping byte sequences to patch embeddings and converting patch-level representations back into byte sequences, respectively~(\autoref{fig:arch_high_level}).

\subsection{Latent Global Transformer Model}

\textit{The Latent Global Transformer} refers to an autoregressive transformer architecture, denoted as $\mathcal{G}$ and consisting of $l_{\mathcal{G}}$ layers, which takes a series of latent input patch encodings $p_j$ and produces a corresponding sequence of output patch encodings $o_j$. In this manuscript, indices $j$ and $i$ are consistently used to represent patches and bytes, respectively. The global transformer employs a block-causal attention mask~\citep{dubey2024llama}, thereby ensuring that attention at any step is limited to patches up to and including the present one within the given document. Since this component dominates the overall computational cost—both in terms of pre-training and inference—strategic selection of its invocation enables flexible allocation of computational resources to various segments of the input and output based on their respective complexities.

\subsection{Local Encoder}
\label{section:local}

\textit{The Local Encoder Model}, symbolized as $\mathcal{E}$, is designed as a compact transformer architecture with significantly fewer layers ($l_{\mathcal{E}} << l_{\mathcal{G}}$) compared to the global encoder. Its chief function is to transform a stream of input bytes $b_i$ into informative patch-level representations $p_j$ with high efficiency. Notably, this model diverges from conventional transformer designs by incorporating a cross-attention layer after each transformer block. This specialized layer aggregates individual byte embeddings into corresponding patch representations, as visualized in~(\autoref{fig:crossattn}).

Initially, each byte in the input sequence, $b_i$, is mapped to an embedding vector using a parameter matrix of dimensions $\mathbb{R}^{256 \times h_{\mathcal{E}}}$, producing the representation $x_i$. Optionally, these embeddings may be enriched with hash-based features as discussed in~(\S\ref{sec:hash-emb}). The model then applies a stack of alternating transformer and cross-attention layers~(\S\ref{sec:enc-cross-attn}), recursively updating the representations until they form patch representations $p_i$, which are subsequently consumed by the global transformer $\mathcal{G}$. Within the transformer modules, \textit{local block causal} attention masks are enforced, allowing each byte to access information within a context window of $w_{\mathcal{E}}$ preceding bytes. This window is permitted to span across different dynamic patches but remains constrained by document boundaries. Detailed information about the embedding mechanism and the structure of the cross-attention module is presented in the following subsections.

\subsubsection{Encoder Hash n-gram Embeddings}
\label{sec:ngram-emb}
\label{sec:hash-emb}
An essential aspect of generating strong and informative representations at every step $i$ lies in embedding contextual data derived from previous bytes. In {\textsc BLT}, this is accomplished by capturing both the identity of the current byte $b_i$ alone and its contribution within a surrounding byte n-gram sequence. At each step $i$, the model initially assembles byte-grams

\begin{align}
    g_{i,n}=\{b_{i-n + 1},\ldots, b_i\}
\end{align}

for each byte position $i$ and $n$ from three to eight.\footnote{
We omit byte-grams of size $n$ or more when $i<n$. }

Next, we present hash $n$-gram embeddings, which assign each byte $n$-gram to a position in a pre-defined embedding table $E_{n}^{hash}$ of constant size using a hash function, with this process performed separately for each $n \in\{3,4,5,6,7,8\}$~\citep{bai2010learning}. The resulting hash $n$-gram embedding is subsequently summed with the embedding corresponding to the individual byte, after which it undergoes normalization and serves as the input to the local encoder. The computation of the combined, or augmented, embedding is as follows:

\begin{align}
e_i &= x_i + \sum_{n = 3,...,8}E_{n}^{hash}(\text{Hash}(g_{i,n})) \\
\text{where, Hash}(g_{i,n}) &= \text{RollPolyHash}(g_{i,n}) \% |E_{n}^{hash}|
\end{align}

We scale $e_i$ by dividing by the total count of $n$-gram sizes incremented by one, employing RollPolyHash as specified in Appendix~\ref{appendix:rollpolyhash}. Section~\ref{section:ablations} presents an analysis of how varying $n$ in $n$-gram hash embeddings, as well as altering the embedding table size, influences the observed trends in flop-controlled scaling laws. Beyond the use of hash-based $n$-gram embeddings, we investigated embeddings derived from $n$-gram frequency information; a comprehensive account of this investigation is given in Appendix~\ref{appendix:ngrams}.

\subsubsection{Encoder Multi-Headed Cross-Attention}
\label{sec:enc-cross-attn}
Our approach is largely inspired by the input cross-attention mechanism described in the Perceiver architecture~\citep{jaegle2021perceiver}. However, a key distinction lies in utilizing latent variables that capture individual patch information, replacing the original model's fixed latent set~(\autoref{fig:crossattn}). Each latent variable is associated with a specific patch and restricts its attention to the bytes within its corresponding patch. In this setup, a dedicated query vector is constructed for every patch $p_j$ by applying a pooling operation over the bytes of patch $p_j$, after which a linear projection, $\mathcal{E}_{C} \in \mathbb{R}^{h_{\mathcal{E}} \times (h_{\mathcal{E}} \times U_{\mathcal{E}})}$, is used. Here, $U_{\mathcal{E}}$ refers to the total number of cross-attention heads in the encoder. More formally, defining $f_{\text{bytes}}(p_j)$ as the ordered list of bytes that comprise patch $p_j$, the process is as follows:

\begin{align}
P_{0,j} &= \mathcal{E}_{C}(f_\text{bytes}((p_j)), f~ \text{is a pooling function} \\
P_l &= P_{l-1} + W_o\left(\text{softmax}\left(\frac{QK^T}{\sqrt{d_k}}\right)V\right) \\
\text{where } Q_j &= W_q(P_{l-1,j}), K_i = W_k(h_{l-1,i}), V_i = W_v(h_{l-1,i}) \\
h_l &= \text{Encoder-Transformer-Layer}_l(h_{l-1})
\end{align}

Here, $P \in \mathbb{R}^{n_p \times h_{\mathcal{G}}}$ denotes the $n_p$ patch embeddings that serve as input to the global model. These embeddings are initialized by aggregating the byte-level representations $e_i$ associated with each patch $p_j$. The matrices $W_q$, $W_k$, $W_v$, and $W_o$ act as linear transformations for the query, key, value, and output spaces, respectively. Notably, both keys and values are formed by projecting the byte-level hidden states $h_i$ from the previous network layer (or $e_i$ if it is the initial layer). A patch-specific masking mechanism ensures that every query $Q_j$ focuses exclusively on keys and values derived from the bytes within its corresponding patch $j$. In the context of multi-headed attention, since patch representations are generally higher dimensional ($h_{\mathcal{G}}$) than $h_{\mathcal{E}}$, we structure $P_l$ as several attention heads, each of dimensionality $h_{\mathcal{E}}$, for cross-attention purposes; the resulting head outputs are concatenated afterward to achieve the $h_{\mathcal{G}}$ size. Furthermore, pre-LayerNorm is applied to queries, keys, and values, while positional encodings are omitted from this cross-attention framework. The module concludes with the addition of a residual connection encircling the cross-attention operation.

\subsection{Local Decoder} 

The local decoder $\mathcal{D}$, like the local encoder, is constructed as an efficient transformer architecture with significantly fewer layers, $l_{\mathcal{D}} << l_{\mathcal{G}}$. Its role is to reconstruct the original byte sequence $y_i$ by decoding a sequence of global patch embeddings $o_j$. To generate the output sequence of bytes, the local decoder utilizes not only the prior decoded bytes but also receives the hidden states from the local encoder corresponding to these bytes. The model comprises $l_{\mathcal{D}}$ stacked modules, each consisting of a cross-attention mechanism followed by a transformer block. In each module, cross-attention is performed first to translate patch-level information into byte-level representations. Subsequently, the transformer block processes these representations as part of the byte sequence.

\subsubsection{Decoder Multi-headed Cross-Attention} 

Within the decoder’s cross-attention mechanism, the functions of queries and keys/values are swapped: here, the byte representations serve as the queries, while the patch representations are used as keys and values. The starting byte representations for this cross-attention step are set to the byte embeddings derived from the output of the final encoder layer, denoted by $h_{l_{\mathcal{E}}}$. For any subsequent layer $l$, the byte representation $d_{l,i}$ is derived as follows:

\begin{align}
D_0 &= h_{l_{\mathcal{E}}} \\
B_l &= D_{l-1} + W_o\left(\text{softmax}\left(\frac{QK^T}{\sqrt{d_k}}\right)V\right), \\
  &\text{where } Q_i = W_q(d_{l-1,i}), K_i = W_k(\mathcal{D}_C(o_j)), V_i = W_v(\mathcal{D}_C(o_j)) \\
  D_l &= \text{Decoder-Transformer-layer}_l(B_l)
\end{align}

Here, $W_k$ and $W_v$ once more serve as projection matrices for the key and value, respectively, each applying a linear transformation and the split operator $\mathcal{D}_C$ to the global model’s final patch outputs $o_j$. The query projection matrix, $W_q$, is applied to byte-level representations $d_{l-1}$ produced by the previous decoder transformer layer (or to $h_{l_{\mathcal{E}}}$ in the initial layer). The output projection, $W_o$, generates the matrix $B \in \mathbb{R}^{h_{\mathcal{D}} \times n_b}$, with $n_b$ denoting the total number of bytes generated as output. Subsequently, the decoder transformer layer computes the following set of decoder representations, $D_l$, using the result $B$ from the cross-attention stage. As with local encoder cross-attention, our implementation uses multi-head attention, applies LayerNorms prior to each sublayer, omits positional encodings, and incorporates a residual connection surrounding the cross-attention.

\section{Experimental Setup}
\label{section:experiments}
We implement a set of systematically controlled experiments to directly compare {\textsc BLT} with models that rely on tokenization, with careful measures to ensure that {\textsc BLT} is not granted any unfair advantages, such as increased sequence context lengths.

\subsection{Pre-training Datasets} In this study, all model sizes utilized are initially trained on two primary datasets: 1) The Llama 2 dataset~\citep{touvron2023llama}, containing 2 trillion tokens aggregated from a broad spectrum of publicly accessible sources. This dataset undergoes rigorous cleaning and filtering processes to enhance data quality; and 2) {\textsc BLT}-1T: A newly constructed corpus comprising 1 trillion tokens sourced from an assortment of public data, which also incorporates a segment of the pre-training material released in Datacomp-LM~\citep{li2024datacomp}. The Llama 2 dataset is specifically employed for scaling law investigations concerning the ideal token count, as established by~\cite{dubey2024llama}, to inform optimal architectural selections for {\textsc BLT}. Conversely, {\textsc BLT}-1T is leveraged for a full-scale pre-training session, aiming to evaluate its performance relative to Llama 3 on subsequent downstream benchmarks. Notably, data derived from Meta products or services are excluded from both datasets. Additionally, when conducting baseline comparisons for models based on tokenizers, we adopt the Llama 3 tokenizer, featuring a vocabulary of 128K tokens, which consistently yielded superior baseline outcomes over the Llama 2 tokenizer during our assessments.

\subsection{Entropy Model} In our study, the entropy model functions as a byte-level language model, which is trained using the identical data distribution as the primary {\textsc BLT} model. Unless specified otherwise, our default entropy model is a transformer comprising 100 million parameters, arranged over 14 layers, with each layer having a hidden size of 512, and a sliding window attention mechanism spanning 512 bytes. All other hyperparameter settings match those utilized in both the local and global transformer variants. Our investigation also included variations in model scale, architectural design, and receptive field, details of which can be found in \autoref{section:ablations}. Notably, for models with sufficiently constrained receptive fields, the trained entropy model lends itself to compact representation via a lookup table.

\subsection{Entropy Threshold and Equalizing Context Length}  
For models that utilize entropy-based patching, we determine a patching threshold designed to produce a specific average \textit{patch size} across the pretraining data mixture. In {\textsc BLT}, as opposed to methods involving tokenization, the \textit{patch size} is a flexible parameter, which notably influences the effective context size accessible to the model. To ensure parity in average context length across models and to prevent those with larger patch sizes from gaining undue benefit, we keep the expected total number of bytes per batch constant. Consequently, when working with models that use larger patch sizes, we proportionally decrease their sequence lengths. Specifically, for Llama 2 data, an 8k byte context is employed, whereas for the {\textsc BLT}-1T dataset, the context window is expanded to 16k bytes on average, all while holding the average batch size fixed at 16M bytes.

Although the mean batch size does not vary, dynamic patching approaches produce varying proportions of bytes per patch when assembling data batches. To optimize efficiency, our implementation of {\textsc BLT} organizes patch batches in a way that eliminates the need for padding within the resource-intensive latent transformer module, resulting in each batch containing an identical patch count. During model training, byte sequences are padded or, if necessary, truncated to fixed lengths of 12k and 24k bytes for the Llama 2 and {\textsc BLT}-1T datasets, respectively, thereby mitigating the risk of memory surges caused by sequences that include exceptionally large patches.

\subsection{Entropy Model Context}
\label{section:entropy-context}
Through experimentation, we observe that applying entropy patching tends to result in increasingly sizable patches when dealing with structured inputs such as multiple choice tasks (refer to patching applied on an MMLU case in \autoref{fig:mmlu_patching}), which are characterized by their repetitive structure. These patterns arise due to decreased entropy in the entropy model context, which typically contains the repeated material. Therefore, for the large-scale application of {\textsc BLT}-Entropy using a patch size of 4.5, we refresh the entropy context by introducing new lines and implement an approximate monotonicity constraint, as this approach is more resistant to "entropy drift" stemming from fluctuations in context length. This adjustment modifies only the manner in which we calculate entropies; our method for determining the appropriate entropy threshold remains unchanged.

\subsection{FLOPs Estimation} 
\label{section:flops}

Our approach to calculating transformer flops is mainly aligned with the methodology in Chinchilla~\citep{hoffmann2022training}, which accounts for the flops involved in feed-forward network operations, qkvo projections within self-attention, as well as both attention and output projection computations. However, our framework differs in that we treat the input embedding layer as a fast lookup operation, not as a dense matrix product, making its contribution to overall flops negligible (0-flops). Consistent with established research, we assume that the backward pass requires double the number of flops as the forward pass.

In order to determine the flops \textit{per byte} for {\textsc BLT} models, we aggregate the flop counts from the local encoder transformer, the global latent transformer, and the local decoder transformer, incorporating as well the contributions from cross-attention modules within both the encoder and decoder. This is formalized as follows:

\begin{align}
    \text{FL}_{\text{{\textsc BLT}}} &= \text{Transf. FL}(h_{\mathcal{G}}, l_{\mathcal{G}}, m=n_{ctx}/n_p,V=0) / n_p\\
   &+\text{Transf. FL}(h_{\mathcal{E}}, l_{\mathcal{E}}, m=w_{\mathcal{E}}, V=0) \\
   &+ \text{Transf. FL}(h_{\mathcal{D}}, l_{\mathcal{D}}, m=w_{\mathcal{D}}, V=256) \\ 
    &+ \text{Cross Attn. FL}(h_{\mathcal{E}}, l_{\mathcal{E}}, m=n_p, r=n_p/k) \times k/n_p \\ 
   &+ \text{Cross Attn. FL}(h_{\mathcal{D}}, l_{\mathcal{D}}, m=k, r=k/n_p) 
\end{align}

Here, $n_{ctx}$ denotes the length of the input sequence measured in bytes, $n_p$ indicates the patch size, $r$ represents the proportion of queries relative to key/values, and $k$ specifies the ratio between the patch dimension and byte dimension—that is, it counts the number of local model segments concatenated to construct a global model representation (as illustrated with $k=2$ in \autoref{fig:crossattn}). $V$ is the vocabulary size relevant to the output projection, utilized exclusively within the local decoder. The effective context length for attention, denoted as $m$, varies depending on whether the module processes the byte-level or patch-level sequence. We adjust the calculation of attention-related FLOPs for each respective module. Detailed formulas for the computation of FLOPs in both Transformer and Cross-Attention settings can be found in Appendix~\ref{section:flops-appendix}.

\subsection{Bits-Per-Byte Estimation} The concept of perplexity is meaningful solely when a specific tokenizer is utilized, since it quantifies the ambiguity associated with each token. In order to facilitate comparison between models operating at the byte and token levels, and in line with earlier studies~\citep{xue2022byt5, yu2023megabyte, wang2024mambabyte}, we adopt Bits-Per-Byte (BPB) as our metric. BPB serves as a tokenizer-agnostic analogue to perplexity. More formally:

\begin{align}
    \text{BPB}(x)&= \frac{\mathcal{L}_{CE}(\pmb{x})}{\ln(2)\cdot n_{\text{bytes}}}
\end{align}

where the uncertainty over the data $\pmb{x}$ as measured by the sum of the cross-entropy loss is normalized by the total number of bytes in $\pmb{x}$ and a constant. 

\subsection{Transformer Architecture Hyperparameters} In configuring the transformer modules of {\textsc BLT}, encompassing both local and global networks, we largely adhere to the framework outlined in Llama~3~\citep{dubey2024llama}. Specifically, we employ the SwiGLU activation~\citep{swiglu} within feed-forward components, and adopt rotary positional embeddings (RoPE)~\citep{RoPe} parameterized with $\theta=500000$~\citep{xiong2024effective}, applying them exclusively to self-attention layers. RMSNorm \citep{RMSNorm} is utilized for layer normalization throughout. For all self-attention modules reliant on standard, fixed attention masks—such as the \textit{block causal} or \textit{fixed-window block causal} variants—we incorporate Flash attention~\citep{dao2022flashattention} and employ a window size of 512 when using fixed-width attention masks. In contrast, given that our cross-attention modules require the use of dynamic, patch-dependent masks, we leverage Flex Attention\footnote{\url{https://pytorch.org/blog/flexattention}} for efficient, fused computations, which substantially accelerates the training process.

\subsection{{\textsc BLT}-Specific Hyperparameters} 

To assess the performance of {\textsc BLT} models, our experimental design encompasses two primary approaches: the investigation of scaling behavior and the assessment on downstream benchmarks. Accordingly, we analyze models across a spectrum of parameter sizes, specifically 400M, 1B, 2B, 4B, and 8B. Details of the architectural hyperparameters employed in these configurations can be found in Appendix~Table~\ref{tab:arch}.

For query initialization in the first cross-attention layer within the local encoder, we employ max-pooling. Consistently across all {\textsc BLT} variants, we utilize $500,000$ hashes paired with a single hash function, and n-gram sizes are chosen between 3 and 8. The learning rate is uniformly set to $4e-4$ for every model. This particular rate was selected after conducting a hyperparameter sweep in the interval $[1e-3, 1e-4]$ at the 400M and 1B scales, which indicated that an identical learning rate yielded optimal results for both token and {\textsc BLT} models.

When analyzing scaling phenomena on the Llama-2 dataset, we adhere to the batch-size recommendations outlined by~\cite{dubey2024llama}, either by sample count or the equivalent byte volume. For optimization, we apply AdamW~\citep{adamw} with hyperparameters $\beta_1 = 0.9$, $\beta_2 = 0.95$, and $\epsilon=10^{-8}$. The learning rate is warmed up linearly over 2000 steps, followed by a cosine decay schedule toward zero. We employ a weight decay coefficient of 0.1 and set a global gradient clipping threshold at 1.0.

\section{Scaling Trends}
\label{section:scaling}

We provide a comprehensive overview of scaling behaviors in byte-level models to guide the continued scaling of {\textsc BLT} architectures. Unlike earlier work on byte-level modeling, our scaling analysis encompasses the following improvements: (a) We evaluate trends under compute-optimal training setups, (b) We train comparable 8B-parameter models using substantial datasets (as large as 1T tokens or 4T bytes) and assess their effectiveness on downstream benchmarks, and (c) We examine scaling characteristics while holding inference costs constant. Subsequent sections will further explore particular benefits that arise from modeling byte sequences.

\subsection{Parameter Matched Compute Optimal Scaling Trends}
\label{section:parameter-matched-scaling}
Leveraging the Llama 2 dataset, we conduct training of multiple \textit{compute-optimal} bpe and {\textsc BLT} models at four different scales, with parameter counts ranging from 1B to 8B. For each configuration, we chart the training flops relative to language modeling accuracy, evaluating performance on a selected subset representative of the entire training mixture. In line with the guidance from Llama~3~\citep{dubey2024llama}, bpe models utilize an empirically derived optimal model-parameter-to-dataset-size ratio. The intention behind this \textit{compute-optimal} methodology is to maximize training dataset performance within a restricted compute allowance~\citep{hoffmann2022training}, establishing a solid comparative point for our study. Each bpe model has an associated {\textsc BLT} variant, which is trained on identical data using a Latent Transformer, carefully configured to mirror the original bpe Transformer’s architecture and scale.

As shown in~\autoref{fig:scaling-compute-opt} (right), {\textsc BLT} models consistently perform as well as or better than their bpe-based counterparts, a pattern that persists with increases in both model size and computation. To our knowledge, {\textsc BLT} represents the first byte-level Transformer architecture to demonstrate scaling behaviors comparable to those of BPE-based models under compute-optimal conditions. These findings support our hypothesis that the parameter-to-compute ratio optimal for bpe models is also applicable to {\textsc BLT}, or, at the very least, closely approximates it.

Achieving bpe scaling trends requires not only advances in architecture but also the implementation of dynamic patching methods. In \autoref{fig:scaling-compute-opt} (left), we present a comparison between models employing space-patching and Llama 3. For this comparison, we approximate SpaceByte \citep{slagle2024spacebyte} by utilizing {\textsc BLT} space-patching while omitting both n-gram embeddings and cross-attention mechanisms. While SpaceByte exhibits performance gains relative to Megabyte, its results still do not approach those of Llama 3. \autoref{fig:scaling-compute-opt} (right) depicts the aggregate benefits arising from both the aforementioned architectural modifications and dynamic patching. At scale, {\textsc BLT} models achieve performance on par with leading tokenizer-based models like Llama 3.

Additionally, we find that the choice of tokenizer has a notable impact on performance for tokenizer-based systems; specifically, models using the Llama-3 tokenizer consistently outperform counterparts trained with the Llama-2 tokenizer when provided with identical training data.

In summary, our {\textsc BLT} architecture demonstrates intermediary behavior between Llama 2 and 3 when larger patch sizes are employed. The average token lengths for the bpe tokenizers in Llama 2 and 3 are 3.7 and 4.4 bytes, respectively. Conversely, {\textsc BLT} can attain comparable scaling performance when configured with patch sizes averaging 6 or even 8 bytes. Since inference flops are inversely related to the mean patch size, utilizing an 8-byte patch size yields up to a 50\% reduction in inference flops. Moreover, models utilizing larger patch sizes appear to benefit more as both model and dataset sizes increase. For example, while {\textsc BLT} with an 8-byte patch size performs markedly worse than bpe Llama 2 at the 1B scale, it surpasses the bpe tokenizer’s performance by the 7B model scale. This observation indicates that such larger patch sizes may be even more advantageous as both model capacity and training compute are increased, and that still greater patch sizes could become practical under continued scaling.

\subsection{Beyond Compute Optimal Task Evaluations}
\label{section:evals}
To further investigate scaling behavior, we train an 8B {\textsc BLT} model past the compute-optimal regime, utilizing the {\textsc BLT}-1T dataset—a more extensive and higher-quality corpus—and examine its performance across a variety of conventional classification and generation benchmarks. For evaluation, we focus on a selection of tasks targeting common sense reasoning, factual knowledge, and code generation:
\paragraph{Classification tasks} encompass ARC-Easy (0-shot)~\citep{Eval_ARC}, ARC-Challenge (0-shot)~\citep{Eval_ARC}, HellaSwag (0-shot)~\citep{Eval_hellaswag}, PIQA (0-shot)~\citep{Eval_piqa}, and MMLU (5-shot)~\citep{hendrycks2020measuring}. We adopt a prompt-based likelihood scoring approach for multiple choice tasks, aggregating the likelihoods assigned to answer choices and reporting mean accuracy as the performance metric.
\paragraph{Coding related generation tasks:} For programming evaluation, we assess models using pass@1 metrics on MBPP (3-shot)~\citep{Eval_MBPP} and HumanEval (0-shot)~\citep{Eval_HumanEval}, providing a measure of the model’s capability to generate syntactically correct and functional Python code.

\autoref{tab:evals} presents a comparison among three models trained on the {\textsc BLT}-1T dataset: a model based on the Llama 3 tokenizer with BPE encoding,\footnote{The Llama 3 tokenizer, which includes a 128k vocabulary, is selected because it outperforms the 32k vocabulary of Llama 2.} and two distinct variants of the {\textsc BLT} architecture. These variants differ in their patching strategies: one leverages a space-patching approach ({\textsc BLT}-Space), while the other applies an entropy-driven patching strategy ({\textsc BLT}-Entropy). The latter operates under an approximate monotonicity constraint and resets the context on new lines, in accordance with the method described in~\autoref{section:entropy-context}. All three models are trained using an equivalent computational (FLOP) budget. Notably, during inference for {\textsc BLT}-Entropy, we adjust the entropy threshold from 0.6 to 0.1, which leads to enhanced task accuracy, albeit requiring additional inference steps.

The {\textsc BLT}-Entropy model achieves superior results compared to the Llama 3 model on 4 of the 7 evaluated tasks, despite both models being trained with an equal number of bytes. This performance advantage likely arises from two factors: (1) more efficient utilization of training computation made possible by dynamic patching, and (2) the explicit modeling of information at the byte level rather than at the token level.

Conversely, {\textsc BLT}-Space yields lower performance compared to the Llama 3 tokenizer in all tasks except one; however, it notably lowers inference FLOPs due to its greater average patch size of 6 bytes. In contrast, both the bpe and entropy-patching methods exhibit an average patch size of about 4.5 bytes within the training data mixture. Given an identical training budget, the model with the larger patch size processes 30\% more data than the other approaches, potentially moving {\textsc BLT} even further from the compute-optimal regime.

\subsection{Patches Scale Better Than Tokens} Utilizing {\textsc BLT} models enables concurrent scaling of both model and patch size without exceeding the designated training and inference flops, and without altering the quantity of training data. This capacity to indefinitely enlarge patch size sets patch-based approaches apart, as they avoid the efficiency constraints intrinsic to fixed-vocabulary token-based methods (see Section~\ref{section:bpe-patch} for further details). Employing longer patches conserves computational resources, which may then be redirected to expand the global latent transformer, since its execution frequency decreases.

In order to evaluate whether increasing model size—while reducing the number of steps on larger patches—yields better performance than utilizing smaller models with more frequent steps, we undertake a fixed inference scaling analysis. We begin with models containing 400 million and 3.6 billion parameters, both employing the Llama 2 tokenizer, and identify models with comparable FLOP counts using the Llama 3 tokenizer. Additionally, we include {\textsc BLT}-Entropy models characterized by average patch sizes of 6 and 8 bytes within the training datamix (refer to~\autoref{tab:fixed-inf-params} for further specifics regarding the models). For those models utilizing a patch size of 8, we implement 3 encoder layers as opposed to just 1. Each configuration is trained across multiple designated training FLOP budgets.

\autoref{fig:fixed_inference_scaling} demonstrates that {\textsc BLT} architectures exhibit superior scaling behavior compared to tokenization-based models across both categories of inference flops. While BPE models initially yield higher performance when trained with limited computational resources, they are soon outperformed by {\textsc BLT} models as the training compute surpasses the regime optimal for the architecture. In practical settings, allocating increased resources during the pretraining phase can result in models that deliver stronger performance under a constrained inference budget. This approach is exemplified by 8B model variants, such as Llama 3.1, which have been exposed to training data volumes nearly two orders of magnitude beyond what is theoretically optimal for their parameter count.

The point at which {\textsc BLT} surpasses token-based architectures has migrated slightly towards the compute-optimal point as models advance to the higher flop class, changing from 3x to 2.5x the compute-optimal budget. Likewise, for the model using a patch size of 8, the rate of scaling is steeper in the larger flop regime, resulting in its performance overtaking the others at an earlier stage. As elaborated in Section~\ref{section:parameter-matched-scaling}, increasing patch size brings performance closer to that of BPE models as the model size grows. This trend can be partly explained by the declining proportion of the total compute that is devoted to the byte-level Encoder and Decoder modules, which do not scale as rapidly as the Latent Transformer. Specifically, when scaling the total number of parameters twenty-fold from 400M to 8B, the localized parameters for {\textsc BLT} expand by only about a factor of two. This nuance is critical, since increases in patch size exclusively influence the computational demand of the patch-level Latent Transformer, leaving the byte-level module computations unaffected. As a result, the {\textsc BLT}-Entropy configuration with ps=8 increased from 1.6x to 1.7x the Llama 2 parameter size when transitioning to a higher model scale.

In conclusion, our investigation into patch-length scaling reveals that the {\textsc BLT} patch-based architecture exhibits superior scaling behavior when both the patch size and model size are expanded concurrently. This positive scaling tendency appears to continue and even strengthen as the scale of the model increases.

\section{Byte Modeling Improves Robustness} Additionally, we evaluate the robustness of {\textsc BLT} relative to token-based models that do not incorporate explicit byte-level signals, and we propose a method for adapting existing token-based models to operate at the byte level.
\label{sec:robustness}

\subsection{Character-Level Tasks}

One initial impetus for developing byte-level models stemmed from their ability to handle noise present at the byte level within inputs, as well as their inherent sensitivity to the substructures comprising tokens—an aspect where models reliant on tokenizers often fall short. To systematically assess these characteristics, we conduct supplementary evaluations on a set of benchmarks designed to test resilience to noisy inputs and sensitivity to character-level features, encompassing both English and multilingual contexts, and including elements such as digits and phonemes. The outcomes of these evaluations are summarized in Table~\ref{tab:char_tasks}.

\paragraph{Noisy Data} To assess the robustness of {\textsc BLT} in comparison to tokenizer-based models, we generate perturbed versions of the benchmark classification datasets detailed in Section~\ref{section:evals}. Five unique strategies are used to introduce character-level noise into the text: (a) \textit{AntSpeak}: The entire text is rendered in uppercase with each character separated by spaces. (b) \textit{Drop}: 10\% of the characters are randomly deleted from the text. (c) \textit{RandomCase}: 50\% of characters are randomly switched to uppercase and the remaining 50\% to lowercase. (d) \textit{Repeat}: 20\% of the characters are randomly selected and each is repeated as many as four consecutive times. (e) \textit{UpperCase}: All characters in the text are converted to uppercase letters. For each evaluation, one of these perturbations is applied to the prompt, the completion, or both independently, and average results are computed across the scenarios. As shown in Table~\ref{tab:char_tasks}, evaluations on the noised version of HellaSwag~\citep{Eval_hellaswag} reveal that {\textsc BLT} consistently demonstrates greater robustness compared to tokenizer-based models, displaying an average improvement of 8 points relative to models trained with identical data, and even surpassing the Llama 3.1 model trained on a significantly larger corpus.

\paragraph{Phonology - Grapheme-to-Phoneme (G2P)} We evaluate the proficiency of {\textsc BLT} in translating a series of graphemes (letters encoding a word) into its phonemic representation (the corresponding pronunciation). The outcomes of the G2P experiment are displayed in Table \ref{tab:char_tasks}, utilizing a 5-shot configuration with Phonology Bench~\citep{suvarna2024phonologybench}. Our analysis shows that {\textsc BLT} achieves superior performance on this task when compared to the Llama 3 1T tokenizer-based baseline model.

\paragraph{CUTE}
To evaluate the character-level comprehension capabilities, we subject {\textsc BLT} to the CUTE benchmark~\citep{edman2024cute}, which consists of multiple tasks grouped into three main types: compositional understanding, recognition of orthographic similarity, and sequence manipulation skills. The CUTE benchmark represents a formidable test for most models that rely on tokenization, as these models tend to retain some knowledge regarding how their tokens are spelled but encounter difficulties when required to use this knowledge for textual manipulations. As presented in Table~\ref{tab:char_tasks}, the {\textsc BLT}-Entropy variant achieves over 25 points higher than either BPE-based Llama 3 model on the CUTE tasks. Our approach stands out particularly in tasks involving manipulation at the character level, achieving near-perfect scores (99.9\%) for both evaluated spelling-related challenges. The magnitude of improvement—despite {\textsc BLT} being trained on only one-sixteenth the data of Llama 3.1—underscores the challenge BPE models face in acquiring fine-grained character-level representations. Figure~\ref{fig:cute_examples} showcases instances in which Llama 3’s tokenizer approach falters while {\textsc BLT} delivers strong performance. The only instances where BPE approaches demonstrate an edge are in word deletion and insertion tasks. This suggests that although manipulating words may be less intuitive for a byte-level approach, the performance gap remains modest, and building up from characters to words may present fewer challenges than decomposing from word-level tokens to characters. Consistency in evaluation procedures was maintained across all tasks, using Huggingface’s original prompts. Notably, BPE models could see improvements through more targeted prompt engineering.

\paragraph{Low Resource Machine Translation} We assess the performance of {\textsc BLT} on both directions of translation—into and out of English—for twenty-one low-resource languages spanning six widely-studied language families, utilizing diverse writing systems included in the FLORES-101 benchmark~\citep{goyal2022flores}. SentencePiece BLEU scores are presented in Table~\ref{tab:flores}. Our findings indicate that {\textsc BLT} surpasses the performance of a model employing the Llama 3 tokenizer, offering a 2-point BLEU improvement for translation into English and a 0.5-point gain for translation from English. For language pairs with extensive prior research, {\textsc BLT} exhibits performance similar to or slightly exceeding that of Llama 3. Importantly, across many language pairs from underrepresented language families, {\textsc BLT} consistently outshines Llama 3, highlighting the robust generalization capabilities of byte-based modeling approaches for handling rare byte sequences.

\subsection{Training {\textsc BLT} from Llama 3}
\label{sec:distillation}
In this study, we investigate a methodology enabling {\textsc BLT} models to capitalize on pre-existing, tokenizer-based pre-trained models for enhanced and accelerated convergence during training. Specifically, we accomplish this by initializing the global transformer parameters of {\textsc BLT} using those from a pre-trained Llama 3.1. In the subsequent finetuning stage, we update the global transformer’s weights utilizing a learning rate set to one-tenth of that applied to the local encoder and decoder components, maintaining the optimal number of steps established for Llama 3. A comparative assessment with a baseline {\textsc BLT} model can be found in Table~\ref{tab:distill}. The results indicate that the {\textsc BLT} model initialized from Llama 3.1 markedly surpasses both the original Llama 3 and the baseline {\textsc BLT} models—where all models were allotted equal computational resources in terms of FLOPs. Notably, even when evaluated against our {\textsc BLT}-Entropy variant (see Table~\ref{tab:evals}), which had been trained on a much larger corpus (1T tokens), the {\textsc BLT} model initialized from Llama 3.1 demonstrates improved performance on the MMLU benchmark. This suggests that the proposed strategy offers a promising avenue for considerably reducing training computational requirements.

This approach can be interpreted as converting models that rely on tokenizers into ones that operate without them, thereby transforming a pre-trained LLaMA 3.1 model into a {\textsc BLT} model. For a thorough evaluation, we report both the baseline LLaMA 3.1 model—pre-trained on 15T tokens—in Table \ref{tab:distill}, and directly compare its results with the corresponding {\textsc BLT} variant derived from LLaMA 3. While the performance of our method is only marginally reduced on MMLU and HumanEval benchmarks, there is a more noticeable decrease across additional tasks. These outcomes highlight the necessity for further exploration to better capitalize on pre-trained models, with a particular focus on refining the data mixtures and tuning additional hyperparameters.

\section{Ablations and Discussion}
\label{section:ablations}
This section presents a series of ablations designed to motivate architectural design decisions for {\textsc BLT}, as well as to evaluate the patching methodology and hyper-parameter selections utilized in the construction and training of the {\textsc BLT} 8B parameter model on the {\textsc BLT}-1T dataset.

\paragraph{Entropy Model Hyper-parameters} We systematically explore the impact of entropy model size and context window length on scaling performance by training a set of byte-level entropy transformer models varying in parameter counts from 1 million to 100 million, and context window lengths ranging between 64 and 512. We produce scaling law plots of bits per byte (bpb) versus training compute (flop), drawing on our $400m$ and $1b$ {\textsc BLT} models trained on the Llama-2 dataset; these results are illustrated in \autoref{fig:entropy_ablation}. The results indicate that improvements in scaling performance are associated with increases in both model size and context window, but gains become marginal once models exceed 50 million parameters.

\paragraph{Types of Patching}

We evaluate the effects of the four patching strategies outlined in Section~\ref{section:patching}, specifically: (1) Strided Patching with strides set to 4 and 6, (2) Patching at positions of whitespace, (3) BPE Tokenizer patching utilizing the Llama 3 tokenizer, and (4) Entropy-based patching leveraging a compact byte-level language model.

Although dynamic patching shortens sequence lengths, we standardize the sequence length across all patching approaches to ensure comparable context windows. Thus, during both training and inference, every model processes the same expected number of bytes per sequence, minimizing confounds associated with broader contextual modeling. The outcomes of these ablation studies are illustrated in Figure~\ref{fig:scaling-compute-opt}. Each of the alternative patching strategies surpasses static patching in performance, with space patching demonstrating results nearly equivalent to those of dynamic entropy-based patching.

In \autoref{tab:evals_ablation}, we report results from standard evaluations of {\textsc BLT} models that utilize three different approaches: tokenizer-based modeling, space patching, and entropy-based patching, each trained with the Llama 2 dataset to the step count considered optimal~\citep{dubey2024llama}. While space patching represents a more straightforward method since it does not require executing an entropy model dynamically during training, our analysis reveals that the improvements gained by entropy-based patching in scaling experiments~(Section~\ref{section:scaling}) persist when applied to downstream benchmark tasks as well.\footnote{The results for space patching are based on earlier experiments without cross-attention; however, comparable patterns have been identified in settings with cross-attention present.}

\paragraph{Cross-Attention} 
In~\autoref{tab:crossattn_ablations_1b}, we present an ablation study examining the effects of integrating cross-attention mechanisms at different locations within the encoder and decoder components of {\textsc BLT}. Regarding the encoder’s cross-attention, we evaluate three methods for query initialization: 1) using a uniform learned embedding across all global states, 2) applying a hash embedding to the patch’s constituent bytes, and 3) aggregating the encoder hidden states corresponding to the bytes within a patch at the respective encoder layer through pooling.

Our results indicate that implementing cross-attention within the \textit{decoder} yields the highest effectiveness. Incorporating cross-attention into the encoder shows marginal gains, and these are primarily observed when queries are initialized via pooling. Furthermore, our analysis demonstrates that cross-attention provides notable advantages on the Common-Crawl dataset, with the benefits being especially pronounced as patch sizes increase.

\paragraph{n-gram Hash Embeddings} 

We evaluate the impact of varying n-gram hash embedding vocabulary sizes at 0, 100K, 200K, and 400K, as summarized in Table~\ref{tab:ngram_ablations_1b}. Our results show that incorporating hash embeddings consistently yields improvements across all tested domains, with particularly notable benefits on Wikipedia and Github (a 0.04 bpb improvement, versus only 0.01 bpb after 15k steps at the 8B scale). At the 8B parameter level, reducing the number of hashes from 500K to 300K leads to an almost negligible performance difference of approximately 0.001 bpb after 15k steps. These findings suggest that hash embeddings play a critical role in aligning the performance of {\textsc BLT} models with that of tokenizer-based baselines. However, beyond 300K hash embeddings, additional gains become marginal. Furthermore, our analysis reveals that improvements from hash embeddings and cross-attention are largely additive, with each offering enhancements on distinct datasets.

\paragraph{Local Model Hyperparamaters} Table \ref{tab:local_layers} presents an ablation of different configurations for the layer counts within the local encoder and decoder components. Our experiments reveal that, in conjunction with hash n-gram embeddings, {\textsc BLT} achieves strong performance when the encoder remains highly minimal—consisting of only a single layer—while the decoder is configured with greater depth.

\section{Related Work}
\label{section:related_work}

\textbf{Character-Level RNNs:} Character language modeling has held longstanding importance since the inception of neural network approaches~\citep{sutskever2011generating,mikolov2012subword,graves2013generating} due to its inherent capacity to natively handle previously unseen words, thereby circumventing the need for back-off strategies. \cite{kim2016character} introduced a system that ingests characters solely on the input side, utilizing convolutional and highway networks integrated with LSTM-based RNNs, achieving parity with leading RNN-based language models for English and demonstrating superior outcomes in morphologically rich languages—a frequently highlighted benefit of character-level LLMs. Expanding on this, \cite{kenter2018byte} deployed byte-level LSTM architectures for machine comprehension, surpassing word-level counterparts on Turkish and Russian, languages with complex morphology. Similarly, \cite{zhang2015character} presented character-level convolutional models for classification, which outperformed word-based models in select scenarios. Further, \citet{chung2019hierarchical} incorporated hierarchical LSTM models coupled with boundary detectors at each tier to uncover underlying textual hierarchies, yielding enhanced character-level language modeling results. In contrast to attention mechanisms, ByteNet~\citep{kalchbrenner2016neural} adopted CNN-based character-level layers for neural machine translation tasks.

\textbf{Character-Level Transformers:} The introduction of transformer architectures leveraging attention mechanisms~\citep{vaswani2017attention} and subword tokenization methods~\citep{sennrich-etal-2016-neural} led to notable advancements in neural language modeling and benchmark evaluations. Nevertheless, the choice of word or sub-word units inherently imposes an inductive bias regarding the granularity at which models process language. Seeking to integrate transformer strengths with early positive outcomes in character-level language modeling, \citet{al2019character} constructed extremely deep transformer networks. By employing auxiliary loss functions, they trained transformer models that surpassed previous character-level LSTM language models. Despite this progress, a marked disparity in performance persisted when compared to word-level large language models. Likewise, GPT-2~\citep{radfordgpt2} demonstrated that, on expansive datasets such as the 1 Billion Word Benchmark, language models operating at the byte-level continued to lag behind their word-level counterparts.

Although \cite{Choe2019BridgingTG} showed that transformer-based LLMs operating at the byte level can surpass subword-level LLMs with similar parameter counts, such byte-level models demand substantially higher computational resources and require more time for training. In a related effort, \cite{el2020characterbert} introduced CharFormer, a BERT variant that constructs word representations via convolutions over character embeddings, yielding better performance in the medical domain, albeit with considerably greater computational expenditure. Furthermore, \cite{clark2022canine} presented CANINE, a 150M-parameter, encoder-only architecture designed to process character sequences directly. This model leverages a deep transformer architecture—reminiscent of the global component in our framework—and combines a local transformer with strided convolutions to downsample input characters. CANINE consistently outperforms mBERT, a comparable token-based encoder-only model, on a range of multilingual downstream tasks. ByT5~\citep{xue2022byt5} explored byte-level encoder-decoder frameworks that avoid all forms of patching or segmentation, yielding models that were more robust to noise and performed on par with traditional tokenizer-based systems even with one quarter the training data. However, the omission of patching operations requires attention computations across all bytes, significantly increasing computational overhead. The transition from subword units to byte-level modeling inherently lengthens input sequences, posing challenges for the scalability of these models. More recently, by employing the Mamba Architecture~\citep{gu2023mamba}, which is capable of maintaining a fixed-size memory state over very extended context windows, \cite{wang2024mambabyte} have developed a byte-level Mamba model without patching that outperforms byte-level transformer baselines on bits-per-byte across multiple datasets at a comparable 350M parameter budget under equalized FLOP constraints.

\textbf{Patching-based approaches:} Utilizing patching strategies offers a means to mitigate the typically high computational cost, measured in flops, associated with byte-level LLMs, all while potentially maintaining competitive performance. Initial promising outcomes at modest scales of both model parameters and training data have been observed in several studies. \cite{nawrot-etal-2022-hierarchical} propose a static patching framework involving downsampling and upsampling modules, leading to the development of the hourglass transformer, which surpasses other byte-level baselines at the 150M parameter scale. \cite{nawrot-etal-2023-efficient} enhance these results through dynamic patching approaches, incorporating mechanisms such as a boundary predictor trained end-to-end, a boundary predictor that leverages supervision from specific tokenizers, and an entropy-driven patching model akin to {\textsc BLT}. Their findings demonstrate that such dynamic patching can yield performance improvements over conventional transformer models of the era when evaluated on language modeling with models containing around 40M parameters and trained on roughly 400M tokens. \cite{lester2024training} examine the impact of training on data sequences compressed via arithmetic coding to achieve greater compression efficiency compared to BPE. By applying an equal-info windows approach, their method surpasses byte-level baselines on language modeling tasks, though it does not exceed the performance of subword-level baselines.

Our research is heavily influenced by, and most directly comparable to, MegaByte~\citep{yu2023megabyte}. MegaByte is a decoder-only, causal large language model that employs a pre-defined, unchanging approach for segmenting bytes into patches through patching and concatenating representations, with a local decoding module that reconstructs the original byte sequence from these patches. Their results show that, at the 1B parameter scale, MegaByte achieves performance competitive with that of models relying on tokenization, evaluated over a corpus of approximately 400B bytes. In our experimental analysis, we systematically examine MegaByte and observe that its static patching strategy does not match the performance of modern, compute-efficient tokenizer-based architectures under fixed computational budgets. We provide evidence that {\textsc BLT} narrows this performance gap. Consistent findings are reported by \cite{slagle2024spacebyte}, who also note the limitations of static patching in MegaByte and propose enhancements, including adapting the patching approach to whitespace and other space-like characters, as well as introducing a local encoder module. Their evaluations reveal better results relative to tokenized transformer models in compute-equivalent experiments within certain datasets, such as Github and arXiv, at the 1B parameter scale. We extend their evaluation in our own study, demonstrating that while these modifications yield progress, further innovations in architecture are essential to truly scale byte-level models and bring their capabilities in line with advanced token-based systems like Llama 3.

\section{Limitations and Future Work}

In this study, our architectural decisions are guided by training models for the empirically optimal number of steps as established for Llama~3~\citep{dubey2024llama}. It is important to note, however, that these scaling relationships were derived for transformers operating at the BPE level, which may not yield the best (data, parameter size) balances when applied to {\textsc BLT}. Determining the scaling laws specific to {\textsc BLT} is deferred to future investigations, and may ultimately reveal even more advantageous scaling behaviors for this architecture. Moreover, most of the present experiments are restricted to a parameter scale of up to 1B, and optimal architectural configurations may shift as we move toward larger model sizes such as 8B parameters and higher, potentially resulting in enhanced performance at greater scales.

Current transformer codebases and libraries are primarily optimized for architectures utilizing tokenizers. Although our work includes theoretical flop-equivalent experiments and leverages some efficient solutions (e.g., FlexAttention) to accommodate layers that diverge from standard transformer structures, there may still be disparities in wall-clock performance when compared to tokenizer-based approaches. Consequently, additional improvements in efficiency could be achieved.

Although {\textsc BLT} relies on an independently trained entropy model for patching, exploring approaches that integrate the patching model into an end-to-end learning framework presents a promising avenue for future research. As demonstrated in Section \ref{sec:distillation}, our preliminary studies reveal encouraging signs when adapting large-scale tokenizer-based models such as Llama 3—trained on over 10T tokens—by transferring their weights to initialize and freeze the global transformer during the "byte-ifying" process. Advancing research in this area may lead to techniques that preserve the advantages of bytefying while potentially exceeding the performance of original tokenizer-based models, all without necessitating training from scratch.

\section{Conclusion}
\label{section:conclusion}

In this work, we introduce the Byte Latent Transformer (\textbf{BLT}), an architectural innovation that challenges the prevailing reliance on fixed-vocabulary tokenization within large language models. {\textsc BLT} employs a dynamic, trainable strategy for aggregating bytes into patches, thereby distributing computation according to the intricacy of input data. This enables notable enhancements in both computational efficiency and resilience to data variability. Our comprehensive scaling analysis reveals that {\textsc BLT} achieves comparable results to prominent token-based architectures such as Llama 3 at scales up to 8B parameters and 4T bytes, and can exchange slight declines in standard evaluation scores for as much as a 50\% reduction in inference flops. Additionally, {\textsc BLT} introduces a novel scaling trajectory, making it possible to jointly expand both model capacity and patch granularity while maintaining a constant inference budget. This flexibility is particularly beneficial under the compute constraints typical of real-world deployment scenarios. By working directly on raw byte sequences, {\textsc BLT} extends model robustness to rare or atypical inputs and enhances the system’s proficiency in modeling sub-word patterns. Collectively, these advancements establish {\textsc BLT} as a compelling substitute for traditional tokenization-driven methodologies, offering an efficient, scalable, and robust platform for future language modeling endeavors.

\end{document}